\documentclass[11pt,a4paper]{article}
\usepackage[utf8]{inputenc}
\usepackage[T1]{fontenc}
\usepackage[french]{babel}
\usepackage[margin=2cm]{geometry}
\usepackage{xcolor}
\usepackage{hyperref}
\usepackage{enumitem}
\usepackage{booktabs}
\usepackage{tabularx}
\usepackage[most]{tcolorbox}

\definecolor{awsblue}{RGB}{35, 47, 62}
\definecolor{awsorange}{RGB}{255, 153, 0}
\definecolor{bgprompt}{RGB}{248, 250, 252}
\definecolor{bgsource}{RGB}{254, 249, 231}

\hypersetup{
    colorlinks=true,
    linkcolor=awsblue,
    urlcolor=blue
}

\newtcolorbox{sourcebox}[1]{
    colback=bgsource,
    colframe=awsorange!85!black,
    fonttitle=\bfseries,
    title={#1},
    breakable,
    enhanced
}

\newtcolorbox{promptbox}[1]{
    colback=bgprompt,
    colframe=awsblue,
    fonttitle=\bfseries,
    title={#1},
    breakable,
    enhanced
}

\title{\textbf{Stratégie Audio \& Podcasts NotebookLM}\\
\Large Certification N\textdegree 1 : AWS Certified Cloud Practitioner (CLF-C02)\\
\large Semaines 1 à 3 (Examen cible : 20 novembre 2026)}
\author{Auren Moyo (Pack Préparation Intensive CDI Data \& AI)}
\date{Novembre 2026}

\begin{document}

\maketitle
\tableofcontents
\newpage

\section{Mode d'Emploi : Générer et Utiliser les 12 Audios CLF-C02}

\subsection{Objectif de la méthode Audio Récursive}
Ce document structure \textbf{12 épisodes de podcasts techniques (débats et discussions approfondies)} à générer sur \textbf{Google NotebookLM} pour couvrir 100 \% du programme officiel \texttt{CLF-C02}.
La progression part des concepts les plus fondamentaux (Niveau 1) jusqu'aux micro-sujets techniques et aux pièges de l'examen (Niveau 3).

\subsection{Procédure étape par étape dans NotebookLM}
Pour chaque épisode (\texttt{CLF-AUDIO-01} à \texttt{CLF-AUDIO-12}) :
\begin{enumerate}[label=\textbf{Étape \arabic* :}]
    \item Crée un nouveau Notebook dans \url{https://notebooklm.google.com/} ou utilise un Notebook dédié \texttt{AWS CLF-C02} en sélectionnant uniquement les sources de l'épisode.
    \item Clique sur \textbf{Add source (Ajouter une source)} $\rightarrow$ \textbf{Discover sources / Recherche Web} et copie-colle les \textbf{Mots-clés de recherche Internet} indiqués dans l'encadré orange.
    \item Ajoute également le ou les \textbf{fichiers Markdown locaux} de ton projet (\texttt{content/certifications/clf/chapters/...}) indiqués dans l'encadré orange.
    \item \textbf{Astuce Pro NotebookLM (Prompt Court + Source Texte) :} Si la boîte \textit{Customize} de l'Audio Overview est limitée à 500 caractères, ajoute d'abord le \textbf{Prompt Complet / Directives détaillées} comme source texte copiée (\textbf{Add source} $\rightarrow$ \textbf{Copied text}, intitulée \texttt{00\_DIRECTIVES\_DEBAT}), puis colle le \textbf{Prompt Court Studio (< 450 caractères)} dans le bouton \textbf{Customize} de l'\textbf{Audio Overview} (langue réglée sur \textbf{Français}).
    \item Une fois l'audio généré, télécharge le fichier et renomme-le avec le nom exact indiqué (exemple : \texttt{CLF-AUDIO-01.mp3}), puis place-le dans le dossier du projet :
    \begin{center}
        \texttt{public/audios/clf-c02/}
    \end{center}
\end{enumerate}

\subsection{Planning d'Écoute Récursive (Semaines 1 à 3)}
Pour ancrer durablement chaque concept dans ta mémoire auditive (transports, marche, révisions du soir), chaque audio suit un cycle de \textbf{3 écoutes espacées} :
\begin{itemize}
    \item \textbf{Écoute 1 (Découverte active, vitesse 1.0x)} : Le jour où tu étudies le bloc correspondant.
    \item \textbf{Écoute 2 (Consolidation, vitesse 1.25x)} : à \textbf{J+2 ou J+3} pour réactiver la mémoire avant l'oubli.
    \item \textbf{Écoute 3 (Rappel rapide \& pièges, vitesse 1.25x ou 1.5x)} : Le \textbf{week-end} et en \textbf{Semaine 3} avant les examens blancs.
\end{itemize}

\begin{table}[h!]
\centering
\small
\begin{tabularx}{\textwidth}{l l X l}
\toprule
\textbf{Code Audio} & \textbf{Type \& Niveau} & \textbf{Thème Principal} & \textbf{Planning d'Écoute} \\
\midrule
\texttt{CLF-AUDIO-01} & Macro (Niv. 1) & Avantages du Cloud, Infra Globale \& 6 Piliers Well-Architected & S1 : Lun, Jeu, Dim \\
\texttt{CLF-AUDIO-02} & Micro (Niv. 1) & AWS Cloud Adoption Framework (CAF) : les 6 perspectives & S1 : Mar, Ven, Dim \\
\texttt{CLF-AUDIO-03} & Micro (Niv. 2) & Modèle de Responsabilité Partagée (IaaS, Managed, Serverless) & S1 : Mer, Sam, S3 \\
\texttt{CLF-AUDIO-04} & Micro (Niv. 2) & IAM en profondeur, MFA, IAM Identity Center, Organizations \& SCP & S1 : Jeu, Dim, S3 \\
\texttt{CLF-AUDIO-05} & Micro (Niv. 3) & Match des Services de Sécurité (GuardDuty, Inspector, Macie, WAF...) & S1 : Ven, S2 Lun, S3 \\
\texttt{CLF-AUDIO-06} & Macro (Niv. 2) & Compute AWS : EC2, Auto Scaling, ELB, Lambda, ECS, EKS, Fargate & S1 : Sam, S2 Mar, S3 \\
\texttt{CLF-AUDIO-07} & Micro (Niv. 3) & Options d'achat EC2 : On-Demand, Reserved, Savings Plans, Spot, Hosts & S2 : Lun, Jeu, S3 \\
\texttt{CLF-AUDIO-08} & Micro (Niv. 2) & Stockage : Classes S3, Lifecycle, EBS vs EFS vs FSx vs Storage Gateway & S2 : Mar, Ven, S3 \\
\texttt{CLF-AUDIO-09} & Micro (Niv. 2) & Bases de Données : RDS (Multi-AZ vs Read Replica), Aurora, DynamoDB... & S2 : Mer, Sam, S3 \\
\texttt{CLF-AUDIO-10} & Micro (Niv. 3) & Réseau VPC : IGW vs NAT, SG vs NACL, Endpoints, CloudFront vs Global Acc. & S2 : Jeu, Dim, S3 \\
\texttt{CLF-AUDIO-11} & Micro (Niv. 2) & Facturation, Cost Explorer vs Budgets vs CUR \& les 5 Plans de Support & S2 : Ven, S3 Lun, S3 Jeu \\
\texttt{CLF-AUDIO-12} & Débat (Niv. 3) & Les 25 Pièges Mortels et Services Homonymes de l'examen CLF-C02 & S2 : Dim, S3 Mar, S3 Jeu \\
\bottomrule
\end{tabularx}
\end{table}

\newpage

% ============================================================================
\section{Partie 1 : Concepts Cloud \& Architecture Globale (Niveau 1)}
% ============================================================================

\subsection{Épisode 01 : Avantages du Cloud, Infrastructure Globale \& Well-Architected}
\begin{sourcebox}{Sources \& Mots-clés Web NotebookLM : \texttt{CLF-AUDIO-01.mp3}}
\textbf{Nom du fichier à enregistrer :} \texttt{CLF-AUDIO-01.mp3}\\
\textbf{Type :} Macro-Thème (Niveau 1 : Fondamental)\\
\textbf{Mots-clés à taper dans la recherche Web NotebookLM (Discover Sources) :}
\begin{itemize}[noitemsep]
    \item \texttt{AWS Six Advantages of Cloud Computing official documentation}
    \item \texttt{AWS Global Infrastructure Regions Availability Zones Local Zones Edge Locations}
    \item \texttt{AWS Well-Architected Framework six pillars official guide CLF-C02}
\end{itemize}
\textbf{Fichiers locaux du projet à importer :}
\texttt{content/certifications/clf/chapters/chapitre_01_concepts_cloud_economie.md} et \texttt{chapitre_02_infrastructure_globale_well_architected.md}
\end{sourcebox}

\begin{promptbox}{Prompt de personnalisation Audio NotebookLM : \texttt{CLF-AUDIO-01}}
\textbf{Prompt Court Studio (< 450 caractères) :}\\
Débat technique en français pour la certification AWS CLF-C02. Décortiquez les 6 avantages officiels du Cloud AWS, la différence entre Haute Disponibilité, Tolérance aux Pannes, Scalabilité et Élasticité, la hiérarchie de l'Infrastructure Globale (Régions, 4 critères de choix, AZ, Edge Locations, Local Zones, Wavelength, Outposts) et les 6 piliers du Well-Architected Framework.

\medskip\hrule\medskip
\textbf{Prompt Complet / Directives détaillées du Débat :}\\
Organisez une discussion technique approfondie et dynamique en français entre deux experts AWS qui préparent un candidat à la certification AWS Certified Cloud Practitioner (CLF-C02). Conservez les termes techniques officiels en anglais.

Votre débat doit décortiquer en détail les 3 blocs suivants et mettre en lumière tous les pièges d'examen :

1. Les 6 avantages officiels du Cloud Computing selon AWS :
Expliquez chacun avec un exemple concret : passer des dépenses en capital (CapEx) aux dépenses variables (OpEx), bénéficier d'économies d'échelle massives, ne plus deviner la capacité nécessaire, augmenter la vitesse et l'agilité, ne plus dépenser d'argent pour gérer et maintenir des datacenters, et se déployer à l'échelle mondiale en quelques minutes.
Piège d'examen à débattre : ne pas confondre la Haute Disponibilité (High Availability), la Tolérance aux Pannes (Fault Tolerance), la Scalabilité (Verticale vs Horizontale) et l'Élasticité.

2. L'Infrastructure Globale AWS :
Décortiquez la hiérarchie exacte : Régions (et les 4 critères de choix d'une Région : conformité légale et gouvernance des données, latence de proximité avec les utilisateurs, disponibilité des services dans la région, et coûts), Availability Zones (AZ : un ou plusieurs datacenters isolés avec alimentation et réseau redondants, interconnectés à faible latence), Points of Presence (Edge Locations et Regional Edge Caches pour CloudFront), AWS Local Zones, AWS Wavelength (5G) et AWS Outposts (hybride sur site).

3. Les 6 Piliers du AWS Well-Architected Framework :
Comparez un par un les 6 piliers : Operational Excellence (opérations sous forme de code, petits changements fréquents et réversibles), Security, Reliability (récupération automatique, scale horizontal), Performance Efficiency, Cost Optimization, et Sustainability (durabilité environnementale). Donnez pour chaque pilier le mot-clé déclencheur dans une question d'examen qui permet d'identifier immédiatement le bon pilier.
\end{promptbox}

\subsection{Épisode 02 : AWS Cloud Adoption Framework (CAF) et ses 6 Perspectives}
\begin{sourcebox}{Sources \& Mots-clés Web NotebookLM : \texttt{CLF-AUDIO-02.mp3}}
\textbf{Nom du fichier à enregistrer :} \texttt{CLF-AUDIO-02.mp3}\\
\textbf{Type :} Micro-Thème (Niveau 1 : Précision sur un sujet très piégeux)\\
\textbf{Mots-clés à taper dans la recherche Web NotebookLM (Discover Sources) :}
\begin{itemize}[noitemsep]
    \item \texttt{AWS Cloud Adoption Framework CAF six perspectives Business People Governance}
    \item \texttt{AWS CAF Platform Security Operations capabilities CLF-C02 exam}
    \item \texttt{AWS 7 Rs migration strategies Rehost Replatform Refactor Retire Retain}
\end{itemize}
\textbf{Fichier local du projet à importer :}
\texttt{content/certifications/clf/chapters/chapitre_12_migration_innovation_caf.md}
\end{sourcebox}

\begin{promptbox}{Prompt de personnalisation Audio NotebookLM : \texttt{CLF-AUDIO-02}}
\textbf{Prompt Court Studio (< 450 caractères) :}\\
Débat ultra-ciblé en français pour la certification AWS CLF-C02 sur le AWS Cloud Adoption Framework (CAF) et les stratégies de migration. Expliquez comment ne jamais confondre les 6 perspectives du CAF (Business, People, Governance, Platform, Security, Operations) avec les 6 piliers Well-Architected, les 4 phases du CAF, et les 7 R de migration.

\medskip\hrule\medskip
\textbf{Prompt Complet / Directives détaillées du Débat :}\\
Animez un débat ultra-ciblé en français dédié à un micro-sujet qui fait perdre énormément de points à l'examen AWS CLF-C02 : le AWS Cloud Adoption Framework (AWS CAF) et les stratégies de migration.

L'un des intervenants doit jouer le rôle d'un examinateur AWS qui pose des mises en situation pièges, et l'autre doit expliquer comment faire la différence immédiate :

1. La distinction fondamentale entre le Well-Architected Framework (comment concevoir techniquement une architecture dans le cloud autour de 6 piliers) et le Cloud Adoption Framework (CAF : comment transformer l'organisation, les équipes et les processus métier pour réussir la migration vers le cloud autour de 6 perspectives).

2. Décortiquez les 6 Perspectives du AWS CAF en les regroupant en 2 familles :
- Les 3 perspectives orientées Capacités Métier (Business Capabilities) :
  * Business (aligner l'IT sur les résultats commerciaux, monétisation des données, gestion de portefeuille).
  * People (culture, leadership transformationnel, formation des équipes, fluidité des compétences cloud).
  * Governance (gestion des risques, gestion financière du cloud, conformité réglementaire).
- Les 3 perspectives orientées Capacités Techniques (Technical Capabilities) :
  * Platform (architecture de la plateforme, ingénierie des données, provisionnement CI/CD).
  * Security (gouvernance de sécurité, gestion des identités, détection des menaces, protection des données).
  * Operations (observabilité, gestion des incidents et des problèmes, continuité d'activité).

3. Les 4 phases de transformation du CAF : Envision (Imaginer), Align (Aligner), Launch (Lancer), Scale (Passer à l'échelle).

4. Les stratégies de migration (les 7 R) : Rehost (Lift and Shift), Replatform (Lift, Tinker and Shift, par exemple migrer une base vers RDS), Repurchase (passer à un SaaS), Refactor / Re-architect (réécrire en cloud-native serverless), Retire, Retain, et Relocate (VMware Cloud on AWS).
\end{promptbox}

\newpage

% ============================================================================
\section{Partie 2 : Sécurité, Identité \& Responsabilité Partagée (Niveau 2 à 3)}
% ============================================================================

\subsection{Épisode 03 : Le Modèle de Responsabilité Partagée dans les moindres détails}
\begin{sourcebox}{Sources \& Mots-clés Web NotebookLM : \texttt{CLF-AUDIO-03.mp3}}
\textbf{Nom du fichier à enregistrer :} \texttt{CLF-AUDIO-03.mp3}\\
\textbf{Type :} Micro-Thème (Niveau 2 : Cœur du Domaine Sécurité 30 \%)\\
\textbf{Mots-clés à taper dans la recherche Web NotebookLM (Discover Sources) :}
\begin{itemize}[noitemsep]
    \item \texttt{AWS Shared Responsibility Model Infrastructure Container Abstracted services}
    \item \texttt{Security OF the Cloud vs Security IN the Cloud AWS CLF-C02}
\end{itemize}
\textbf{Fichier local du projet à importer :}
\texttt{content/certifications/clf/chapters/chapitre_03_modele_responsabilite_iam.md}
\end{sourcebox}

\begin{promptbox}{Prompt de personnalisation Audio NotebookLM : \texttt{CLF-AUDIO-03}}
\textbf{Prompt Court Studio (< 450 caractères) :}\\
Débat contradictoire en français sur le Modèle de Responsabilité Partagée AWS pour l'examen CLF-C02. Détaillez la frontière exacte des responsabilités entre AWS et le client selon les 3 catégories de services : IaaS (EC2 : patching OS à la charge du client), Services Gérés (RDS : patching OS et backups par AWS) et Services Abstraits (S3, DynamoDB, Lambda).

\medskip\hrule\medskip
\textbf{Prompt Complet / Directives détaillées du Débat :}\\
Organisez un débat contradictoire en français sur le Modèle de Responsabilité Partagée d'AWS (Shared Responsibility Model) pour la certification CLF-C02.

Au lieu de rester sur la théorie générale (AWS est responsable de la sécurité DU Cloud, le client est responsable de la sécurité DANS le Cloud), entrez dans le détail chirurgical des 3 catégories de services et des pièges d'examen :

1. Ce qui est TOUJOURS à 100 \% sous la responsabilité d'AWS :
Sécurité physique des datacenters, destruction des disques physiques en fin de vie, câblage réseau physique, serveurs hôtes physiques et couche de virtualisation (Hyperviseur Nitro/Xen).

2. Ce qui est TOUJOURS à 100 \% sous la responsabilité du Client :
Les données du client (classification, intégrité), la gestion des identités et des accès (IAM, mots de passe, MFA), le chiffrement côté client et l'activation du chiffrement côté serveur, et l'authentification des utilisateurs.

3. Les contrôles partagés (Shared Controls) : comment la frontière se déplace selon le type de service :
- Services IaaS (Amazon EC2) : le client gère l'OS invité (correctifs/patching Windows ou Linux), le pare-feu réseau (Security Groups), les applications installées et le chiffrement des volumes EBS.
- Services Conteneurisés / Gérés (Amazon RDS, Amazon ElasticBeanstalk, EKS) : sur RDS, AWS prend en charge le patching de l'OS et les mises à jour mineures du moteur de base de données ainsi que les sauvegardes automatiques, mais le client reste responsable du schéma de base de données, des utilisateurs SQL, des Security Groups et de l'activation du chiffrement KMS.
- Services Abstraits / Serverless (Amazon S3, DynamoDB, AWS Lambda) : AWS gère l'infrastructure, l'OS et la plateforme ; le client configure le code (Lambda), les permissions IAM/Bucket Policies et le chiffrement des données.

4. Testez-vous mutuellement sur 5 cas pièges d'examen (qui met à jour le firmware matériel ? qui patche l'OS d'une instance EC2 ? qui patche l'OS d'une base RDS ? qui configure les règles d'un Security Group ?).
\end{promptbox}

\subsection{Épisode 04 : IAM en profondeur, MFA, IAM Identity Center \& AWS Organizations}
\begin{sourcebox}{Sources \& Mots-clés Web NotebookLM : \texttt{CLF-AUDIO-04.mp3}}
\textbf{Nom du fichier à enregistrer :} \texttt{CLF-AUDIO-04.mp3}\\
\textbf{Type :} Micro-Thème (Niveau 2 : Sécurité des Identités)\\
\textbf{Mots-clés à taper dans la recherche Web NotebookLM (Discover Sources) :}
\begin{itemize}[noitemsep]
    \item \texttt{AWS IAM Users Groups Roles Policies Least Privilege Root User best practices}
    \item \texttt{AWS IAM Identity Center vs IAM Roles vs AWS Organizations SCP}
\end{itemize}
\textbf{Fichier local du projet à importer :}
\texttt{content/certifications/clf/chapters/chapitre_03_modele_responsabilite_iam.md}
\end{sourcebox}

\begin{promptbox}{Prompt de personnalisation Audio NotebookLM : \texttt{CLF-AUDIO-04}}
\textbf{Prompt Court Studio (< 450 caractères) :}\\
Débat technique en français pour l'examen CLF-C02 sur la gestion des identités AWS : protection du Root User, IAM Users, Groups, Roles (identifiants temporaires STS sans Access Keys en dur), évaluation des IAM Policies JSON (Explicit Deny gagne toujours), et distinction entre AWS IAM Identity Center (SSO) et AWS Organizations (Service Control Policies SCP).

\medskip\hrule\medskip
\textbf{Prompt Complet / Directives détaillées du Débat :}\\
Créez un podcast technique en français décortiquant toute l'architecture d'identité et de contrôle d'accès sur AWS pour l'examen CLF-C02.

Abordez point par point avec rigueur :
1. Le compte Root AWS : pourquoi il ne faut jamais créer d'Access Keys pour le Root User, l'activation obligatoire du MFA dès la création du compte, et les rares tâches qui nécessitent exclusivement le Root User (fermer le compte AWS, changer de plan de support AWS, modifier les informations de facturation ou le nom du compte).
2. IAM Users, IAM Groups et IAM Policies :
- Un IAM Group ne peut contenir QUE des IAM Users (jamais d'autres groupes, jamais d'instances EC2).
- Anatomie d'une politique IAM JSON : Effect (Allow ou Deny), Action, Resource, Condition, Principal.
- La règle d'or d'évaluation IAM : tout est refusé par défaut (Implicit Deny), un Allow explicite autorise l'action, mais un Deny explicite (Explicit Deny) l'emporte TOUJOURS sur n'importe quel Allow.
3. IAM Roles et AWS STS (Security Token Service) :
Pourquoi une application tournant sur une instance EC2 ou une fonction Lambda ne doit JAMAIS contenir d'Access Keys en dur dans le code, mais doit utiliser un IAM Role (Instance Profile) qui fournit des identifiants temporaires via STS.
4. Les moyens d'accéder à AWS : Console Management (mot de passe + MFA), AWS CLI et AWS SDKs (Access Key ID + Secret Access Key), et AWS CloudShell (terminal pré-authentifié dans le navigateur).
5. La différence cruciale entre :
- AWS IAM Identity Center (anciennement AWS SSO) : authentification unique centralisée pour les employés humains à travers plusieurs comptes AWS et applications SaaS.
- AWS Organizations et les Service Control Policies (SCP) : les SCP définissent le plafond maximal d'autorisations (garde-fous) pour une Organizational Unit (OU) ou un compte membre (y compris le Root User du compte membre), mais une SCP n'accorde jamais d'autorisation à elle seule.
\end{promptbox}

\subsection{Épisode 05 : Le Grand Match des Services de Sécurité \& Conformité AWS}
\begin{sourcebox}{Sources \& Mots-clés Web NotebookLM : \texttt{CLF-AUDIO-05.mp3}}
\textbf{Nom du fichier à enregistrer :} \texttt{CLF-AUDIO-05.mp3}\\
\textbf{Type :} Micro-Thème (Niveau 3 : Piège n\textdegree 1 de la certification CLF-C02)\\
\textbf{Mots-clés à taper dans la recherche Web NotebookLM (Discover Sources) :}
\begin{itemize}[noitemsep]
    \item \texttt{AWS GuardDuty vs Inspector vs Macie vs Detective vs Security Hub}
    \item \texttt{AWS Shield Standard Advanced vs AWS WAF vs Firewall Manager}
    \item \texttt{AWS KMS vs CloudHSM vs Secrets Manager vs Systems Manager Parameter Store}
    \item \texttt{AWS Artifact vs Audit Manager vs Trusted Advisor security checks}
\end{itemize}
\textbf{Fichier local du projet à importer :}
\texttt{content/certifications/clf/chapters/chapitre_04_securite_conformite_avances.md}
\end{sourcebox}

\begin{promptbox}{Prompt de personnalisation Audio NotebookLM : \texttt{CLF-AUDIO-05}}
\textbf{Prompt Court Studio (< 450 caractères) :}\\
Débat en français sous forme de duels comparatifs sur les services de sécurité AWS pour l'examen CLF-C02 : GuardDuty (menaces ML sur logs) vs Inspector (vulnérabilités CVE sur EC2/ECR/Lambda) vs Macie (données sensibles PII sur S3) vs Detective vs Security Hub ; Shield vs WAF vs Firewall Manager ; KMS vs CloudHSM ; Secrets Manager vs Parameter Store ; Artifact vs Audit Manager.

\medskip\hrule\medskip
\textbf{Prompt Complet / Directives détaillées du Débat :}\\
Animez un débat ultra-rythmé en français consacré au piège numéro 1 de l'examen AWS CLF-C02 : la confusion entre les services de sécurité, de chiffrement et de conformité d'AWS.

Construisez la discussion sous forme de duels comparatifs avec les mots-clés exacts qui tombent dans les énoncés d'examen :

Duel 1 : Détection des menaces vs Vulnérabilités vs Données sensibles :
- Amazon GuardDuty : détection intelligente des menaces (Machine Learning) qui analyse en continu les logs CloudTrail, VPC Flow Logs et DNS Logs (ex: minage de cryptomonnaie, exfiltration, comportement suspect).
- Amazon Inspector : scan automatisé des vulnérabilités logicielles (CVE) et des expositions réseau non intentionnelles sur les instances EC2, les images de conteneurs dans ECR et le code des fonctions Lambda.
- Amazon Macie : découverte et protection des données personnelles (PII, numéros de cartes bancaires, passeports) stockées exclusivement dans les buckets Amazon S3 grâce au Machine Learning.
- Amazon Detective : investigation forensique et recherche de la cause racine (Root Cause Analysis) après une alerte GuardDuty.
- AWS Security Hub : tableau de bord centralisé qui agrège les alertes et vérifie la conformité par rapport aux standards (CIS, PCI-DSS).

Duel 2 : Protection Réseau et Applicative :
- AWS Shield Standard (gratuit et automatique contre les DDoS couches 3 et 4) vs AWS Shield Advanced (payant, protection sophistiquée, équipe SRT 24/7, protection financière contre les pics de coûts liés à un DDoS).
- AWS WAF (Web Application Firewall) : couche 7 (HTTP/HTTPS), protège contre les injections SQL et le Cross-Site Scripting (XSS) devant CloudFront, ALB ou API Gateway.
- AWS Firewall Manager : gestion centralisée des règles WAF, Shield et Security Groups à travers toute une AWS Organization.

Duel 3 : Chiffrement et Gestion des Secrets :
- AWS KMS (Key Management Service : clés gérées mutualisées) vs AWS CloudHSM (boîtier matériel dédié FIPS 140-2 Level 3 à locataire unique).
- AWS Secrets Manager (rotation automatique native des mots de passe de bases de données RDS) vs AWS Systems Manager Parameter Store (stockage hiérarchique de paramètres et configurations).
- AWS Certificate Manager (ACM) : certificats SSL/TLS gratuits et renouvelés automatiquement.

Duel 4 : Conformité :
- AWS Artifact : portail en libre-service pour télécharger les rapports d'audit et certificats de conformité d'AWS (ISO, SOC, PCI) et gérer les accords BAA (HIPAA).
- AWS Audit Manager : collecte continue et automatisée de preuves pour préparer les audits internes du client.
\end{promptbox}

\newpage

% ============================================================================
\section{Partie 3 : Compute, Stockage, Bases de Données \& Réseau (Niveau 2 à 3)}
% ============================================================================

\subsection{Épisode 06 : Services de Calcul (Compute) : EC2, Conteneurs et Serverless}
\begin{sourcebox}{Sources \& Mots-clés Web NotebookLM : \texttt{CLF-AUDIO-06.mp3}}
\textbf{Nom du fichier à enregistrer :} \texttt{CLF-AUDIO-06.mp3}\\
\textbf{Type :} Macro-Thème (Niveau 2 : Cœur du Domaine Technologie 34 \%)\\
\textbf{Mots-clés à taper dans la recherche Web NotebookLM (Discover Sources) :}
\begin{itemize}[noitemsep]
    \item \texttt{AWS EC2 instance types Auto Scaling Elastic Load Balancing CLF-C02}
    \item \texttt{AWS Lambda vs ECS vs EKS vs Fargate vs Elastic Beanstalk vs Lightsail vs Batch}
\end{itemize}
\textbf{Fichier local du projet à importer :}
\texttt{content/certifications/clf/chapters/chapitre_05_services_calcul_ec2_lambda_conteneurs.md}
\end{sourcebox}

\begin{promptbox}{Prompt de personnalisation Audio NotebookLM : \texttt{CLF-AUDIO-06}}
\textbf{Prompt Court Studio (< 450 caractères) :}\\
Débat technique en français pour la certification CLF-C02 sur les services Compute AWS : familles d'instances EC2, Elastic Load Balancing (ALB, NLB, GWLB) et Auto Scaling, conteneurs (ECS vs EKS vs ECR et le moteur serverless AWS Fargate), fonctions serverless AWS Lambda (limite 15 min), Elastic Beanstalk (PaaS), Lightsail et AWS Batch.

\medskip\hrule\medskip
\textbf{Prompt Complet / Directives détaillées du Débat :}\\
Organisez une discussion pédagogique et comparative en français sur tous les services de Calcul (Compute) d'AWS pour la certification CLF-C02.

Décortiquez précisément :
1. Amazon EC2 et ses familles d'instances :
Comment retenir facilement les familles : General Purpose (M, T : équilibré, serveurs web), Compute Optimized (C : calcul intensif, batch, encodage vidéo), Memory Optimized (R, X : bases de données in-memory, SAP HANA), Accelerated Computing (P, G, Inf : GPU et Machine Learning), Storage Optimized (I, D : IOPS locales très élevées, bases NoSQL).
2. Haute Disponibilité et Scalabilité avec EC2 :
Différence entre Elastic Load Balancing (ALB pour HTTP/HTTPS, NLB pour TCP/UDP ultra-rapide, GWLB pour les appliances de sécurité) et EC2 Auto Scaling (qui ajoute ou supprime automatiquement des instances selon la demande).
3. Conteneurs sur AWS :
- Amazon ECS (orchestrateur de conteneurs Docker natif AWS) vs Amazon EKS (Kubernetes géré sur AWS) vs Amazon ECR (registre privé d'images de conteneurs).
- Le rôle clé d'AWS Fargate : moteur de calcul serverless pour conteneurs (fonctionne avec ECS et EKS) qui évite d'avoir à provisionner ou gérer des instances EC2 sous-jacentes.
4. Serverless et Plateformes gérées (PaaS) :
- AWS Lambda : fonctions événementielles serverless, facturation à la milliseconde et au nombre de requêtes, durée d'exécution maximale stricte de 15 minutes.
- AWS Elastic Beanstalk : PaaS où le développeur fournit uniquement son code (Java, Python, Node.js) et Beanstalk provisionne EC2, Auto Scaling, ELB et RDS (le service Beanstalk est gratuit, on ne paie que les ressources sous-jacentes).
- Amazon Lightsail : serveurs privés virtuels (VPS) à prix mensuel fixe et prévisible pour des sites simples (WordPress) sans expertise cloud.
- AWS Batch : exécution de milliers de jobs de calcul batch à grande échelle (utilise des instances EC2 ou Fargate).
\end{promptbox}

\subsection{Épisode 07 : Les 7 Options d'Achat EC2 (Micro-Thème Critique)}
\begin{sourcebox}{Sources \& Mots-clés Web NotebookLM : \texttt{CLF-AUDIO-07.mp3}}
\textbf{Nom du fichier à enregistrer :} \texttt{CLF-AUDIO-07.mp3}\\
\textbf{Type :} Micro-Thème (Niveau 3 : Arbitrage financier et technique très fréquent)\\
\textbf{Mots-clés à taper dans la recherche Web NotebookLM (Discover Sources) :}
\begin{itemize}[noitemsep]
    \item \texttt{AWS EC2 purchasing options On-Demand Reserved Savings Plans Spot}
    \item \texttt{AWS Dedicated Hosts vs Dedicated Instances vs Capacity Reservations CLF-C02}
\end{itemize}
\textbf{Fichiers locaux du projet à importer :}
\texttt{content/certifications/clf/chapters/chapitre_05_services_calcul_ec2_lambda_conteneurs.md} et \texttt{chapitre_11_facturation_tarification_support.md}
\end{sourcebox}

\begin{promptbox}{Prompt de personnalisation Audio NotebookLM : \texttt{CLF-AUDIO-07}}
\textbf{Prompt Court Studio (< 450 caractères) :}\\
Débat ultra-ciblé en français sur les 7 options d'achat EC2 pour l'examen CLF-C02 : On-Demand, Spot Instances (-90 \%, préavis d'interruption de 2 min), Compute Savings Plans (inclut Fargate et Lambda) vs EC2 Instance Savings Plans, Standard vs Convertible Reserved Instances, et le piège Dedicated Hosts (visibilité sockets physiques, licences BYOL) vs Dedicated Instances.

\medskip\hrule\medskip
\textbf{Prompt Complet / Directives détaillées du Débat :}\\
Animez un débat ultra-précis en français consacré exclusivement aux 7 options d'achat et de réservation d'Amazon EC2. C'est l'un des sujets les plus testés à l'examen CLF-C02 : le candidat doit savoir choisir la bonne option selon l'énoncé.

Comparez en détail chaque option, son pourcentage de réduction, son engagement et son cas d'usage :
1. On-Demand Instances : facturation à la seconde (Linux) ou à l'heure, aucun engagement initial, idéal pour les charges imprévisibles ou les nouveaux projets à court terme.
2. Spot Instances : jusqu'à 90 \% de réduction par rapport au tarif On-Demand, mais AWS peut interrompre l'instance avec un préavis de 2 minutes si la capacité est réclamée. Idéal pour les traitements batch tolérants aux pannes, l'analyse de données et les jobs qui peuvent reprendre sur point de sauvegarde. Interdit pour les bases de données critiques !
3. Savings Plans (engagement d'un montant en dollars par heure sur 1 an ou 3 ans) :
- Compute Savings Plans : le plus flexible (jusqu'à 66 \% d'économie), s'applique automatiquement à EC2 (quelle que soit la région, la famille, la taille ou l'OS), mais aussi à AWS Fargate et AWS Lambda !
- EC2 Instance Savings Plans : jusqu'à 72 \% d'économie, lié à une famille d'instances dans une région donnée (permet de changer la taille ou l'OS).
4. Reserved Instances (RI : engagement sur 1 an ou 3 ans pour une capacité spécifique) :
- Standard RI (jusqu'à 72 \% d'économie, revente possible sur la RI Marketplace) vs Convertible RI (jusqu'à 66 \% d'économie, permet d'échanger la famille d'instance, l'OS ou la location).
5. Dedicated Hosts vs Dedicated Instances (le piège classique) :
- Dedicated Hosts : serveur physique entièrement dédié qui donne la visibilité sur les sockets physiques et les cœurs matériels. Indispensable pour les licences logicielles liées au matériel (BYOL : Bring Your Own License, Windows Server, SQL Server, Oracle) et la conformité stricte.
- Dedicated Instances : instances qui tournent sur du matériel dédié à un seul client, mais sans visibilité ni contrôle sur les sockets physiques ni l'affinité d'hôte lors d'un redémarrage.
6. On-Demand Capacity Reservations : permet de réserver de la capacité EC2 dans une Availability Zone précise pour la durée souhaitée sans engagement de 1 ou 3 ans.
\end{promptbox}

\subsection{Épisode 08 : Stockage AWS : Classes Amazon S3, EBS, EFS, FSx \& Storage Gateway}
\begin{sourcebox}{Sources \& Mots-clés Web NotebookLM : \texttt{CLF-AUDIO-08.mp3}}
\textbf{Nom du fichier à enregistrer :} \texttt{CLF-AUDIO-08.mp3}\\
\textbf{Type :} Micro-Thème (Niveau 2 : Stockage Objet, Bloc, Fichier et Hybride)\\
\textbf{Mots-clés à taper dans la recherche Web NotebookLM (Discover Sources) :}
\begin{itemize}[noitemsep]
    \item \texttt{Amazon S3 Storage Classes Standard Intelligent-Tiering One Zone-IA Glacier}
    \item \texttt{AWS EBS vs EFS vs Instance Store vs Amazon FSx vs Storage Gateway}
\end{itemize}
\textbf{Fichier local du projet à importer :}
\texttt{content/certifications/clf/chapters/chapitre_06_services_stockage_s3_ebs_efs.md}
\end{sourcebox}

\begin{promptbox}{Prompt de personnalisation Audio NotebookLM : \texttt{CLF-AUDIO-08}}
\textbf{Prompt Court Studio (< 450 caractères) :}\\
Débat technique en français pour la certification CLF-C02 sur le stockage AWS : comparaison des 7 classes de stockage Amazon S3 (Standard, Intelligent-Tiering, Standard-IA, One Zone-IA, Glacier Instant/Flexible/Deep Archive), S3 Object Lock (WORM), et différence entre Instance Store, EBS (1 AZ), EFS (multi-AZ Linux NFS), FSx (Windows/Lustre) et Storage Gateway.

\medskip\hrule\medskip
\textbf{Prompt Complet / Directives détaillées du Débat :}\\
Organisez un débat technique en français sur le stockage AWS pour l'examen CLF-C02.

Décortiquez deux grands axes :
Axe 1 : Amazon S3 (Stockage Objet, durabilité de 99,999999999 \% soit 11 neufs) et ses Classes de Stockage :
- S3 Standard : données fréquemment consultées (répliquées sur au moins 3 AZ).
- S3 Intelligent-Tiering : déplace automatiquement les objets entre les niveaux d'accès fréquent et peu fréquent grâce au monitoring des accès, sans frais de récupération (idéal quand les schémas d'accès sont inconnus ou changeants).
- S3 Standard-IA (Infrequent Access) : données peu consultées mais nécessitant un accès immédiat en millisecondes en cas de besoin (au moins 3 AZ, frais de récupération par Go).
- S3 One Zone-IA : stocké dans UNE SEULE Availability Zone (20 \% moins cher que Standard-IA), idéal pour les copies secondaires de sauvegarde ou les données facilement recréables (comme des miniatures d'images). Perdu si l'AZ est détruite !
- S3 Glacier Instant Retrieval (archives consultées 1 fois par trimestre avec accès en millisecondes), S3 Glacier Flexible Retrieval (récupération de 1 minute à 12 heures), et S3 Glacier Deep Archive (le moins cher d'AWS, conservation 7 à 10 ans, récupération en 12 à 48 heures).
- Mentionnez S3 Lifecycle Policies, S3 Versioning et S3 Object Lock (WORM : Write Once Read Many).

Axe 2 : Comparatif Stockage Bloc vs Fichier vs Hybride :
- Instance Store : stockage bloc éphémère physiquement attaché à l'hôte (données perdues si l'instance EC2 est arrêtée).
- Amazon EBS (Elastic Block Store) : disque virtuel attaché à une instance EC2 dans une MÊME AZ, snapshots sauvegardés sur S3.
- Amazon EFS (Elastic File System) : système de fichiers réseau partagé Linux (NFS) accessible simultanément par des centaines d'instances EC2 à travers plusieurs AZ.
- Amazon FSx : systèmes de fichiers tiers gérés (FSx for Windows File Server en SMB/Active Directory, FSx for Lustre pour le HPC et le Machine Learning haute performance).
- AWS Storage Gateway : pont hybride reliant le datacenter on-premises au stockage cloud S3 (File Gateway, Volume Gateway, Tape Gateway).
\end{promptbox}

\subsection{Épisode 09 : Bases de Données \& Services Analytics AWS}
\begin{sourcebox}{Sources \& Mots-clés Web NotebookLM : \texttt{CLF-AUDIO-09.mp3}}
\textbf{Nom du fichier à enregistrer :} \texttt{CLF-AUDIO-09.mp3}\\
\textbf{Type :} Micro-Thème (Niveau 2 : Bases relationnelles, NoSQL et Data Analytics)\\
\textbf{Mots-clés à taper dans la recherche Web NotebookLM (Discover Sources) :}
\begin{itemize}[noitemsep]
    \item \texttt{Amazon RDS Multi-AZ vs Read Replicas Amazon Aurora DynamoDB Redshift}
    \item \texttt{AWS Analytics Athena Glue Kinesis EMR QuickSight OpenSearch CLF-C02}
\end{itemize}
\textbf{Fichier local du projet à importer :}
\texttt{content/certifications/clf/chapters/chapitre_07_bases_de_donnees_analytics.md}
\end{sourcebox}

\begin{promptbox}{Prompt de personnalisation Audio NotebookLM : \texttt{CLF-AUDIO-09}}
\textbf{Prompt Court Studio (< 450 caractères) :}\\
Débat en français pour l'examen CLF-C02 sur les bases de données et l'analytics AWS : RDS Multi-AZ (haute disponibilité synchrone) vs RDS Read Replicas (scalabilité en lecture asynchrone), Amazon Aurora (6 copies sur 3 AZ), DynamoDB (NoSQL serverless), ElastiCache, Neptune (graphes), DocumentDB, et la chaîne Redshift, Athena, Glue, Kinesis, EMR et QuickSight.

\medskip\hrule\medskip
\textbf{Prompt Complet / Directives détaillées du Débat :}\\
Créez un débat structuré en français pour maîtriser parfaitement tous les services de Bases de Données et d'Analytics demandés à la certification AWS CLF-C02.

1. Bases de données relationnelles (SQL) :
- Amazon RDS : moteurs supportés (PostgreSQL, MySQL, MariaDB, Oracle, SQL Server, Db2).
- Le piège classique : RDS Multi-AZ (Haute Disponibilité et bascule automatique synchrone vers une instance standby en cas de panne, mais la standby ne sert PAS à lire des données) vs RDS Read Replicas (jusqu'à 15 réplicas en lecture asynchrone pour soulager la base principale et améliorer la scalabilité des requêtes de reporting).
- Amazon Aurora : base relationnelle cloud-native compatible MySQL et PostgreSQL, 5 fois plus rapide que MySQL, 3 fois plus rapide que PostgreSQL, qui réplique automatiquement 6 copies des données à travers 3 Availability Zones.

2. Bases de données NoSQL et spécialisées :
- Amazon DynamoDB : base NoSQL clé-valeur et documents entièrement serverless, latence inférieure à 10 millisecondes à n'importe quelle échelle, DAX (cache in-memory).
- Amazon ElastiCache : cache en mémoire géré (Redis et Memcached) pour accélérer les bases de données.
- Amazon DocumentDB (compatible MongoDB), Amazon Keyspaces (compatible Apache Cassandra), Amazon Neptune (base de données orientée Graphes : réseaux sociaux, détection de fraude), Amazon MemoryDB.

3. L'écosystème Data & Analytics AWS :
- Amazon Redshift : entrepôt de données (Data Warehouse) colonnaire pour l'analytique OLAP et la BI.
- AWS Glue : service d'intégration de données et d'ETL serverless avec son Data Catalog.
- Amazon Athena : service de requêtage SQL interactif serverless directement sur les fichiers stockés dans Amazon S3 (paiement à la requête / au volume scanné).
- Amazon Kinesis : ingestion et traitement de flux de données en temps réel (streaming).
- Amazon EMR : traitement Big Data géré (Apache Spark, Hadoop).
- Amazon QuickSight : tableaux de bord de Business Intelligence (BI) serverless avec Machine Learning intégré.
\end{promptbox}

\subsection{Épisode 10 : Réseau VPC, Sécurité Réseau \& Distribution Mondiale}
\begin{sourcebox}{Sources \& Mots-clés Web NotebookLM : \texttt{CLF-AUDIO-10.mp3}}
\textbf{Nom du fichier à enregistrer :} \texttt{CLF-AUDIO-10.mp3}\\
\textbf{Type :} Micro-Thème (Niveau 3 : Architecture Réseau et Pièges VPC)\\
\textbf{Mots-clés à taper dans la recherche Web NotebookLM (Discover Sources) :}
\begin{itemize}[noitemsep]
    \item \texttt{AWS VPC Subnets Internet Gateway NAT Gateway Security Groups vs NACL}
    \item \texttt{AWS VPC Peering vs Transit Gateway vs VPC Endpoints PrivateLink}
    \item \texttt{Amazon CloudFront vs AWS Global Accelerator Route 53 Direct Connect VPN}
\end{itemize}
\textbf{Fichier local du projet à importer :}
\texttt{content/certifications/clf/chapters/chapitre_08_reseau_vpc_cloudfront_route53.md}
\end{sourcebox}

\begin{promptbox}{Prompt de personnalisation Audio NotebookLM : \texttt{CLF-AUDIO-10}}
\textbf{Prompt Court Studio (< 450 caractères) :}\\
Débat technique en français pour la certification CLF-C02 sur le réseau AWS : sous-réseaux publics/privés (Internet Gateway vs NAT Gateway), Security Groups (Stateful, instance, Allow uniquement) vs Network ACLs (Stateless, subnet, Allow et Deny), VPC Endpoints, Site-to-Site VPN vs Direct Connect, Route 53, et CloudFront vs AWS Global Accelerator.

\medskip\hrule\medskip
\textbf{Prompt Complet / Directives détaillées du Débat :}\\
Animez un débat technique en français sur l'architecture Réseau AWS (Amazon VPC et services mondiaux) pour la certification CLF-C02.

Expliquez avec une clarté chirurgicale les comparatifs suivants qui tombent systématiquement à l'examen :
1. Sous-réseau Public vs Sous-réseau Privé dans un VPC :
- Internet Gateway (IGW) : permet aux ressources d'un sous-réseau public de communiquer dans les deux sens avec Internet.
- NAT Gateway : permet aux instances situées dans un sous-réseau privé de sortir vers Internet (pour télécharger des mises à jour logicielles) tout en empêchant Internet d'initier une connexion entrante vers ces instances privées.

2. Security Groups vs Network ACLs (NACL) :
- Security Group : agit au niveau de l'instance (ENI), il est Stateful (si le trafic entrant est autorisé, le trafic de réponse sortant est automatiquement autorisé), et il ne contient QUE des règles Allow (pas de règles Deny).
- Network ACL (NACL) : agit au niveau du sous-réseau (Subnet), il est Stateless (il faut autoriser explicitement l'entrée ET la sortie), et il supporte à la fois des règles Allow et des règles Deny évaluées par ordre de numéro.

3. Connectivité Privée et Hybride :
- VPC Peering (connexion 1-à-1 non transitive entre deux VPC) vs AWS Transit Gateway (architecture en étoile Hub-and-Spoke pour interconnecter des centaines de VPC et réseaux on-premises).
- VPC Endpoints (Gateway Endpoints pour S3 et DynamoDB, Interface Endpoints avec PrivateLink) : accéder aux services AWS sans jamais passer par l'Internet public.
- AWS Site-to-Site VPN (connexion chiffrée rapide à mettre en place via l'Internet public) vs AWS Direct Connect (liaison fibre optique physique dédiée et privée entre le datacenter et AWS : latence constante et bande passante garantie, mais prend plusieurs semaines à installer).

4. Services Mondiaux :
- Amazon Route 53 : service DNS mondial hautement disponible.
- Amazon CloudFront (CDN qui met en cache le contenu statique et dynamique dans les Edge Locations) vs AWS Global Accelerator (utilise deux adresses IP statiques Anycast et fait entrer le trafic TCP/UDP immédiatement sur le réseau dorsal privé d'AWS sans mise en cache).
\end{promptbox}

\newpage

% ============================================================================
\section{Partie 4 : Facturation, Support \& les 25 Pièges Ultimes (Niveau 3)}
% ============================================================================

\subsection{Épisode 11 : Facturation, Outils de Coûts \& les 5 Plans de Support AWS}
\begin{sourcebox}{Sources \& Mots-clés Web NotebookLM : \texttt{CLF-AUDIO-11.mp3}}
\textbf{Nom du fichier à enregistrer :} \texttt{CLF-AUDIO-11.mp3}\\
\textbf{Type :} Micro-Thème (Niveau 2 : Domaine Billing, Pricing \& Support 12 \%)\\
\textbf{Mots-clés à taper dans la recherche Web NotebookLM (Discover Sources) :}
\begin{itemize}[noitemsep]
    \item \texttt{AWS Pricing Calculator vs Cost Explorer vs AWS Budgets vs Cost and Usage Report}
    \item \texttt{AWS Support Plans Basic Developer Business Enterprise On-Ramp Enterprise TAM}
\end{itemize}
\textbf{Fichier local du projet à importer :}
\texttt{content/certifications/clf/chapters/chapitre_11_facturation_tarification_support.md}
\end{sourcebox}

\begin{promptbox}{Prompt de personnalisation Audio NotebookLM : \texttt{CLF-AUDIO-11}}
\textbf{Prompt Court Studio (< 450 caractères) :}\\
Débat précis en français pour la certification CLF-C02 sur les outils financiers AWS (Pricing Calculator avant déploiement vs Cost Explorer pour l'analyse et les prévisions à 12 mois vs AWS Budgets pour les alertes de seuil vs Cost and Usage Report CUR), le Consolidated Billing, et le comparatif des 5 plans de Support AWS (Basic, Developer, Business, Enterprise On-Ramp, Enterprise).

\medskip\hrule\medskip
\textbf{Prompt Complet / Directives détaillées du Débat :}\\
Organisez un débat précis en français consacré au Domaine 4 de la certification CLF-C02 : Facturation, Gestion des Coûts et Plans de Support AWS.

1. Les 4 Outils Financiers AWS (ne jamais les confondre à l'examen) :
- AWS Pricing Calculator : estimer le coût d'une nouvelle architecture AVANT de la déployer dans AWS.
- AWS Cost Explorer : visualiser, filtrer et analyser les coûts passés des 12 derniers mois et obtenir une prévision (forecast) des coûts sur les 12 prochains mois, ainsi que des recommandations d'achat de Savings Plans.
- AWS Budgets : définir des seuils de dépenses ou d'utilisation personnalisés et envoyer des alertes proactives (Email, SNS) lorsque le coût réel ou prévu dépasse le seuil.
- AWS Cost and Usage Report (CUR) : le rapport le plus granulaire et détaillé (heure par heure, ligne par ligne avec les Cost Allocation Tags), livré sous forme de fichiers dans un bucket Amazon S3 (analysable avec Athena et QuickSight).
- Consolidated Billing dans AWS Organizations : une seule facture unique pour tous les comptes membres, partage des remises sur volume (Volume Discounts) et partage des Reserved Instances / Savings Plans entre les comptes.

2. Les 5 Plans de Support AWS (comparaison stricte) :
- Basic (gratuit) : Service Health Dashboard, documentation, forums, et 7 contrôles de base sur AWS Trusted Advisor.
- Developer : pour les environnements de test, support par email aux heures ouvrées, 1 contact principal.
- Business : minimum requis pour les charges de production ! Support technique 24/7 par téléphone/chat/email, accès complet à TOUS les contrôles AWS Trusted Advisor, accès à l'API AWS Support, temps de réponse < 1 heure pour un système de production en panne.
- Enterprise On-Ramp : pont intermédiaire pour les entreprises en croissance, accès à un pool partagé de Technical Account Managers (TAM), temps de réponse < 30 minutes pour les charges critiques.
- Enterprise : un Technical Account Manager (TAM) DÉDIÉ à votre entreprise, accès à l'équipe Concierge Support pour la facturation, Infrastructure Event Management (IEM) inclus, et temps de réponse < 15 minutes pour les charges critiques en panne.
\end{promptbox}

\subsection{Épisode 12 : Grand Débat Récapitulatif : Les 25 Pièges et Services Homonymes}
\begin{sourcebox}{Sources \& Mots-clés Web NotebookLM : \texttt{CLF-AUDIO-12.mp3}}
\textbf{Nom du fichier à enregistrer :} \texttt{CLF-AUDIO-12.mp3}\\
\textbf{Type :} Grand Débat Final (Niveau 3 : Révision intensive veille d'examen)\\
\textbf{Mots-clés à taper dans la recherche Web NotebookLM (Discover Sources) :}
\begin{itemize}[noitemsep]
    \item \texttt{AWS CloudWatch vs CloudTrail vs AWS Config vs Systems Manager vs X-Ray}
    \item \texttt{AWS SNS vs SQS vs EventBridge vs Step Functions}
    \item \texttt{AWS AI services Rekognition Textract Comprehend Polly Transcribe Kendra Bedrock}
\end{itemize}
\textbf{Fichiers locaux du projet à importer :}
\texttt{content/certifications/clf/chapters/chapitre_09_integration_decouplage_monitoring.md}, \texttt{chapitre_10_intelligence_artificielle_ml.md} et les 4 fiches dans \texttt{content/certifications/clf/fiches_express/}
\end{sourcebox}

\begin{promptbox}{Prompt de personnalisation Audio NotebookLM : \texttt{CLF-AUDIO-12}}
\textbf{Prompt Court Studio (< 450 caractères) :}\\
Grand épisode final de révision intensive en français pour l'examen AWS CLF-C02 : comparatif rapide entre CloudWatch (métriques/logs) vs CloudTrail (qui a appelé quelle API) vs AWS Config (historique de configuration) vs X-Ray ; SQS (file Pull) vs SNS (Pub/Sub Push) vs EventBridge ; CloudFormation vs SSM ; et tous les services IA (Bedrock, SageMaker, Rekognition, Textract, Comprehend, Kendra, Lex).

\medskip\hrule\medskip
\textbf{Prompt Complet / Directives détaillées du Débat :}\\
Animez un grand épisode final de révision intensive en français pour l'examen AWS CLF-C02, conçu pour être écouté la veille des examens blancs et le matin de l'examen officiel.

Passez en revue sous forme de quiz rapide et d'explications percutantes les paires et trios de services homonymes ou proches :

1. Le Trio Observabilité et Audit (qui tombe à chaque examen) :
- Amazon CloudWatch : performance et monitoring opérationnel (métriques CPU, logs applicatifs, alarmes). Question clé : "Comment se comporte ma ressource ?"
- AWS CloudTrail : audit et traçabilité des appels d'API et actions des utilisateurs. Question clé : "QUI a fait QUOI, QUAND et depuis quelle adresse IP ?"
- AWS Config : inventaire et historique de la configuration des ressources et vérification de leur conformité dans le temps. Question clé : "Quelle était la configuration de ce Security Group mardi dernier ?"
- AWS X-Ray : traçage distribué et débogage de microservices et fonctions Lambda.

2. Le Trio Découplage et Messagerie :
- Amazon SQS (Simple Queue Service) : file d'attente de messages (modèle Pull / Polling) pour découpler deux microservices.
- Amazon SNS (Simple Notification Service) : publication/abonnement (Pub/Sub, modèle Push) pour envoyer une notification à de multiples abonnés (Email, SMS, Lambda, files SQS).
- Amazon EventBridge : bus d'événements serverless permettant de router des événements AWS ou SaaS et de déclencher des tâches planifiées (cron).

3. Infrastructure as Code et Gestion de parc :
- AWS CloudFormation (déployer toute une infrastructure AWS via un template YAML ou JSON) vs AWS CDK (définir l'infrastructure en vrai code Python/TypeScript).
- AWS Systems Manager (SSM) : gérer et patcher à grande échelle une flotte d'instances EC2 et de serveurs on-premises (Session Manager permet de se connecter sans ouvrir le port SSH 22 ni gérer de clés SSH).

4. Toute la famille Intelligence Artificielle et Machine Learning d'AWS :
- Amazon Bedrock (IA générative serverless et Foundation Models via API) vs Amazon Q Developer (assistant IA de codage) vs Amazon SageMaker (plateforme complète pour les Data Scientists qui construisent, entraînent et déploient leurs propres modèles ML).
- Services IA prêts à l'emploi : Amazon Rekognition (images et vidéos), Amazon Textract (extraction OCR de texte, formulaires et tableaux dans des documents scannés), Amazon Comprehend (NLP, analyse de sentiments), Amazon Transcribe (Speech-to-Text), Amazon Polly (Text-to-Speech), Amazon Translate (traduction), Amazon Kendra (moteur de recherche d'entreprise intelligent), Amazon Lex (moteur conversationnel des chatbots derrière Amazon Connect).
\end{promptbox}

\end{document}
